\documentclass[a4paper,10pt]{article}
\usepackage[utf8]{inputenc}
\usepackage[english]{babel}
\usepackage{geometry}
\geometry{margin=1in}
\usepackage{booktabs}
\usepackage{xcolor}
\usepackage{hyperref}

\title{6502 Monitor - ESP32 Pin Mapping Reference}
\author{Atari ESP32 Project}
\date{\today}

\begin{document}

\maketitle

\section{Overview}
This document details the final pin assignment for the ESP32 NodeMCU DevKit V1 interfacing with a 6502 (Atari) system at 1.79 MHz. The firmware utilizes direct register access and IRAM-resident tasks to meet timing requirements.

\section{Data Bus (Bidirectional)}
All data bus pins are located in \textbf{GPIO Bank 0} for optimized 8-bit parallel access.

\begin{table}[h!]
\centering
\begin{tabular}{@{}llll@{}}
\toprule
\textbf{6502 Signal} & \textbf{ESP32 GPIO} & \textbf{Board Label} & \textbf{Level Shifter} \\ \midrule
D0                   & GPIO 4             & D4                  & Yes (5V $\leftrightarrow$ 3.3V) \\
D1                   & GPIO 5             & D5                  & Yes \\
D2                   & GPIO 13            & D13                 & Yes \\
D3                   & GPIO 14            & D14                 & Yes \\
D4                   & GPIO 16            & RX2                 & Yes \\
D5                   & GPIO 17            & TX2                 & Yes \\
D6                   & GPIO 18            & D18                 & Yes \\
D7                   & GPIO 19            & D19                 & Yes \\ \bottomrule
\end{tabular}
\end{table}

\section{Address Bus (Inputs)}
A0--A3 are connected to "Input Only" pins for maximum safety.

\begin{table}[h!]
\centering
\begin{tabular}{@{}llll@{}}
\toprule
\textbf{6502 Signal} & \textbf{ESP32 GPIO} & \textbf{Board Label} & \textbf{Characteristics} \\ \midrule
A0                   & GPIO 34            & D34                 & \textbf{Input Only}      \\
A1                   & GPIO 35            & D35                 & \textbf{Input Only}      \\
A2                   & GPIO 36            & VP                  & \textbf{Input Only}      \\
A3                   & GPIO 39            & VN                  & \textbf{Input Only}      \\
A4                   & GPIO 32            & D32                 & Input/Output             \\
A5                   & GPIO 33            & D33                 & Input/Output             \\
A6                   & GPIO 21            & D21                 & Input/Output             \\
A7                   & GPIO 27            & D27                 & Input/Output             \\ \bottomrule
\end{tabular}
\end{table}

\section{Control Signals (Inputs)}
\begin{table}[h!]
\centering
\begin{tabular}{@{}llll@{}}
\toprule
\textbf{6502 Signal} & \textbf{ESP32 GPIO} & \textbf{Board Label} & \textbf{Description} \\ \midrule
PHI2                 & GPIO 2             & D2                  & System Clock (Synchronization) \\
R/W                  & GPIO 15            & D15                 & Read (H) / Write (L) \\
SEL\_N               & GPIO 22            & D22                 & D1XX\_N (PBI) or CCTL\_N (Cart) \\
ROMSEL               & GPIO 23            & D23                 & \$D800--\$DFFF Range Select (PBI) \\
RAMSEL               & GPIO 26            & D26                 & \$D600--\$D7FF Range Select (PBI) \\ \bottomrule
\end{tabular}
\end{table}

\newpage
\section{Control Signals (Outputs)}
\begin{table}[h!]
\centering
\begin{tabular}{@{}llll@{}}
\toprule
\textbf{6502 Signal} & \textbf{ESP32 GPIO} & \textbf{Board Label} & \textbf{Description} \\ \midrule
EXTSEL               & GPIO 3             & \textbf{RX0}        & Disable Atari Memory (Active Low) \\
VCS                  & GPIO 25            & D25                 & Device Select (Active Low) \\
MPD                  & GPIO 12            & \textbf{D12}        & Math Pack Disable (Active Low) \\ \bottomrule
\end{tabular}
\end{table}

\section{System \& Debug}
\begin{table}[h!]
\centering
\begin{tabular}{@{}llll@{}}
\toprule
\textbf{Signal}      & \textbf{ESP32 GPIO} & \textbf{Board Label} & \textbf{Description} \\ \midrule
Debug TX             & GPIO 1             & TX0                 & Serial output (115200 bps) \\
GND                  & GND                & GND                 & Common Ground with Atari \\ \bottomrule
\end{tabular}
\end{table}

\section{Critical Hardware Notes}
\begin{itemize}
    \item \textbf{GPIO 12 (MPD) Safety:} This is a bootstrap pin. It \textbf{MUST be LOW} during ESP32 power-on/reset. If the Atari system pulls this pin HIGH at boot, the ESP32 will fail to start. Ensure a pull-down resistor or that the level shifter is inactive during ESP32 reset.
    \item \textbf{RX0 (GPIO 3) Usage:} UART0 Receive is disabled in software. This pin is strictly used as a digital output for \textbf{EXTSEL}. Do not connect a serial source (PC RX) to this pin if the Atari is connected.
    \item \textbf{Level Shifters:} All signals must pass through 3.3V $\leftrightarrow$ 5V level shifters (like TXS0108E).
    \item \textbf{PBI vs CCTL Mode:} Toggle the mode in \texttt{main.cpp} via \texttt{\#define BUS_MODE}.
\end{itemize}

\end{document}
